%=====================================================================
% The Yang--Mills Mass Gap
% Converted from standalone paper: yang_mills_mass_gap.tex
%=====================================================================
\section{The Yang--Mills Mass Gap}
\label{ch:yang_mills}

The Yang--Mills existence and mass gap problem is one of the seven Clay Millennium Problems. We prove that the mass gap is a theorem of discrete geometry: it exists on the natural simplicial complex over $\CP^4$, and \emph{cannot} exist on $\RR^4$. The mass gap and continuous spacetime are mutually exclusive.

%=====================================================================
\subsection{Introduction}
\label{sec:intro}
%=====================================================================

The Yang--Mills existence and mass gap problem, one of the seven Clay
Millennium Problems \cite{Clay}, asks:

\begin{quote}
\emph{For any compact simple gauge group $G$, prove that the quantum
Yang--Mills theory on $\RR^4$ has a mass gap $\Delta > 0$.}
\end{quote}

\noindent
Lattice QCD \cite{Wilson, Creutz} provides strong numerical evidence:
Monte Carlo simulations on finite lattices consistently produce a mass
gap, and the result survives extrapolation to the continuum limit. But
numerical computation is not mathematical proof. Decades of effort have
produced no rigorous proof within the continuum framework of the Wightman
\cite{Wightman} or Osterwalder--Schrader \cite{OS} axioms.

We present a resolution from a different direction. In the DRLT framework
, spacetime is not a continuum approximated by a
lattice; spacetime \emph{is} a finite simplicial complex on
$\CCP^4$. The mass gap is then a theorem of discrete algebraic
geometry, following from three facts:

\begin{enumerate}
\item The strong force lives on AAA hinges --- triangles with all three
  vertices in the spatial sector $\CC^3$.
\item Confinement is the combinatorial fact $\binom{3}{3} = 1$: only one
  AAA configuration exists per local patch, exhausting all color degrees
  of freedom in a single hop.
\item The Regge action of a single AAA hinge is strictly positive
  ($\sqrt{\det} \times \delta > 0$), giving a mass gap bounded away
  from zero.
\end{enumerate}

We then prove a no-go theorem: taking the continuum limit
($N \to \infty$, lattice spacing $a \to 0$) forces
$\det(G_{\mathrm{AAA}}) \to 0$, destroying the mass gap. The mass gap
exists if and only if spacetime is discrete.

The chapter is organized as follows.
Section~\ref{sec:framework} reviews the DRLT framework.
Section~\ref{sec:definitions} gives rigorous definitions of the
transfer matrix, mass gap, correlation length, and propagation.
Section~\ref{sec:confinement} proves confinement as $\CC^3$ rank
saturation.  Section~\ref{sec:massgap} proves the mass gap.
Section~\ref{sec:nogo} proves the no-go theorem for continuous
spacetime.  Section~\ref{sec:navier_stokes} formalises the
Navier--Stokes regularity theorem on the lattice and proves that
the continuum limit destroys regularity by the same mechanism.
Section~\ref{sec:discussion} discusses implications.

%=====================================================================
\subsection{Framework}
\label{sec:framework}
%=====================================================================

\subsubsection{The axiom and its unique substrate}

DRLT begins from one axiom: \emph{$N$ points exist with pairwise
relations.} The relations take values in an algebra $\mathbb{K}$.
Four requirements uniquely select $\mathbb{K} = \CC$
\cite{Frobenius}:

\begin{enumerate}
\item[\textbf{R1.}] \textbf{Magnitude} (geometry): $|G_{ij}|$ is
  well-defined. Requires $\mathbb{K}$ normed.
\item[\textbf{R2.}] \textbf{Phase} (gauge forces): $\arg(G_{ij})$ is
  continuous. Requires unit-norm elements to form a connected Lie group.
\item[\textbf{R3.}] \textbf{Determinant} (gravity): $\det(G)$ is a
  single real number. Requires $\mathbb{K}$ commutative.
\item[\textbf{R4.}] \textbf{Winding} (charge quantization): holonomy
  admits discrete winding numbers. Requires
  $\pi_1(\text{phase group}) \cong \mathbb{Z}$.
\end{enumerate}

\noindent
By Frobenius, the finite-dimensional associative division algebras over
$\RR$ are $\{\RR, \CC, \mathbb{H}\}$. $\RR$
fails R2 ($\{+1,-1\}$ is discrete). $\mathbb{H}$ fails R3
(non-commutative). $\CC$ satisfies all four, with
$\pi_1(S^1) = \mathbb{Z}$.

Each point $i$ carries a state $\psi_i \in \CC^d$, and the
pairwise relation is the Gram matrix:
\begin{equation}
G_{ij} = \langle \psi_i | \psi_j \rangle, \qquad
G \in M_{N \times N}(\CC), \qquad
\mathrm{Tr}(G) = N.
\end{equation}

\subsubsection{The $(3,2)$ split and gauge group}

The unique decomposition of $\CC^d$ into atomic blocks supporting
chiral fermions, CP violation, and charge quantization is
$(n_S, n_T) = (3, 2)$, yielding $d = 5$ 
(cf.\ \cite{GeorgiGlashow} for the original SU(5) GUT). This gives:

\begin{itemize}
\item Base space: $\CCP^4$ (projectivization of
  $\CC^5 \setminus \{0\}$).
\item Gauge group: $\mathrm{SU}(3) \times \mathrm{SU}(2) \times
  \mathrm{U}(1)$.
\item Spatial sector: $\CC^3$ (indices $a = 0, 1, 2$).
\item Temporal sector: $\CC^2$ (indices $\alpha = 3, 4$),
  orthogonal to $\CC^3$.
\end{itemize}

The spatial Gram matrix is the restriction:
\begin{equation}
G^S_{ij} = \sum_{a=0}^{2} \psi^*_{i,a}\, \psi_{j,a}, \qquad
\mathrm{rank}(G^S) \leq n_S = 3.
\label{eq:GS}
\end{equation}

\subsubsection{Simplicial complex and Regge calculus}

The $N$ vertices of $G$ induce a finite simplicial complex $K$ on
$\CCP^4$. Each 4-simplex $\sigma$ has 5 vertices; each 2-face
(triangle) $h$ is a \emph{hinge}. The hinge carries:

\begin{itemize}
\item \textbf{Area}: $A_h = \sqrt{\det(G|_h)}$, where $G|_h$ is the
  $3 \times 3$ Gram sub-matrix of the hinge's three vertices.
\item \textbf{Deficit angle}: $\delta_h = 2\pi -
  \sum_{\sigma \supset h} \theta_\sigma(h)$, where $\theta_\sigma(h)$
  is the dihedral angle at $h$ within $\sigma$.
\end{itemize}

\noindent
The Regge action \cite{Regge} is:
\begin{equation}
\boxed{S_R = \sum_{h \in K} A_h\, \delta_h.}
\label{eq:Regge_YM}
\end{equation}

Hinges are classified by the $(3,2)$ split of their vertices:

\begin{center}
\begin{tabular}{clcl}
\hline
Type & Vertices & Count per simplex & Force \\
\hline
AAA & 3 spatial & $\binom{3}{3}\binom{2}{0} = 1$ & Strong \\
AAB & 2 spatial $+$ 1 temporal & $\binom{3}{2}\binom{2}{1} = 6$ &
  Electromagnetic \\
ABB & 1 spatial $+$ 2 temporal & $\binom{3}{1}\binom{2}{2} = 3$ &
  Weak \\
BBB & 3 temporal & $\binom{2}{3} = 0$ & (forbidden) \\
\hline
& & Total: $10 = \binom{5}{3}$ & \\
\hline
\end{tabular}
\end{center}

\noindent
The strong force resides exclusively on the unique AAA hinge per simplex.

%=====================================================================
\subsection{Rigorous Definitions}
\label{sec:definitions}
%=====================================================================

\begin{definition}[Transfer matrix]
\label{def:transfer}
Let $K$ be a simplicial complex on $\CCP^4$ with $N$ vertices.
Define the \emph{transfer matrix} between adjacent spatial slices
$\Sigma_t$ and $\Sigma_{t+1}$ as:
\begin{equation}
T_{mn} = \sum_{\sigma \supset (m,n)}
  \exp\Bigl(-\sum_{h \subset \sigma} A_h\,\delta_h / \hbar_h\Bigr),
\end{equation}
where the sum runs over 4-simplices $\sigma$ connecting vertex $m$
in $\Sigma_t$ to vertex $n$ in $\Sigma_{t+1}$, and
$\hbar_h = A_h/(4\ln 2)$ is the dynamical Planck constant at hinge $h$.
\end{definition}

\begin{definition}[Lattice Hamiltonian and mass gap]
\label{def:massgap}
The lattice Hamiltonian $H$ is defined through $T = e^{-aH}$, where
$a$ is the lattice spacing.  Let $E_0 < E_1 \leq E_2 \leq \cdots$ be
the eigenvalues of $H$ in the gauge-singlet sector.  The
\emph{mass gap} is:
\begin{equation}
\Delta \;=\; E_1 - E_0.
\end{equation}
\end{definition}

\begin{definition}[Correlation length]
\label{def:correlation}
The \emph{correlation length} $\xi$ in the strong sector is defined by
the exponential decay of the connected two-point function:
\begin{equation}
\langle \mathcal{O}(x)\, \mathcal{O}^\dagger(y) \rangle_{\mathrm{conn}}
\;\sim\; e^{-|x - y|/\xi} \qquad \text{as } |x-y| \to \infty,
\end{equation}
where $\mathcal{O}$ is a gauge-invariant operator in the
$\mathrm{SU}(3)$ sector and distance is measured in lattice units.
The mass gap and correlation length are related by $\Delta = 1/\xi$.
\end{definition}

\begin{definition}[Propagation]
\label{def:propagation}
A gauge excitation \emph{propagates} from hinge $h_1$ to hinge $h_2$
if there exists a chain of adjacent hinges $h_1, h_1', \ldots, h_2$
of the same type such that $\langle \mathcal{O}(h_1)\,
\mathcal{O}^\dagger(h_2) \rangle_{\mathrm{conn}} > 0$.
The \emph{propagation range} is the maximum distance (in hinge hops)
over which this correlator is non-zero.
\end{definition}

%=====================================================================
\subsection{Confinement as $\CC^3$ Rank Saturation}
\label{sec:confinement}
%=====================================================================

\begin{definition}[Effective channel count]
For a hinge type with $k$ spatial and $(3-k)$ temporal vertices, the
number of independent configurations is:
\begin{equation}
N_{\mathrm{eff}}(k) = \binom{n_S}{k} \binom{n_T}{3-k}.
\end{equation}
\end{definition}

\noindent
The channel counts for each force are:

\begin{center}
\begin{tabular}{lccc}
\hline
Force & Hinge type & $N_{\mathrm{eff}}$ & Propagation range \\
\hline
Strong & AAA & $\binom{3}{3}\binom{2}{0} = 1$ & 1 hop (confined) \\
Electromagnetic & AAB & $\binom{3}{2}\binom{2}{1} = 6$ & $\infty$ \\
Weak & ABB & $\binom{3}{1}\binom{2}{2} = 3$ & 3 hops (short-range) \\
\hline
\end{tabular}
\end{center}

\begin{proposition}[Confinement]
\label{prop:confinement}
The strong force cannot propagate beyond one hinge. Any gauge-invariant
excitation in the $\mathrm{SU}(3)$ sector is confined to a single AAA
triangle.
\end{proposition}

\begin{proof}
Propagation from hinge $h_1$ to an adjacent hinge $h_2$ requires
consuming an independent channel at each hop. The first hop uses the
unique AAA configuration $(S_1, S_2, S_3)$, which spans all of
$\CC^3$. A second hop requires a \emph{second independent} AAA
configuration --- three new linearly independent spatial vectors. But
$\mathrm{rank}(G^S) \leq n_S = 3$, and the first triple already
saturates this rank. Any further spatial vector is a linear combination
of the first three, producing $\det(G|_{h_2}) = 0$: a degenerate hinge.
Propagation is algebraically forbidden.
\end{proof}

\begin{remark}
This is a combinatorial fact, not a dynamical phenomenon. Confinement
in DRLT requires no running coupling, no flux tubes, no Wilson loops,
and no dynamical symmetry breaking. It is the statement that
$\CC^3$ has exactly three dimensions.
\end{remark}

\begin{remark}
The correlation length in the strong sector is exactly $\xi = 1$ (one
lattice unit). The mass gap, defined as the inverse correlation length,
is therefore $\Delta = 1/\xi > 0$ in lattice units. The remainder of
this paper computes the \emph{value} of this gap from the Regge action.
\end{remark}

%---------------------------------------------------------------------
\subsubsection{Confinement as solution non-existence}
\label{sec:confinement_solution}
%---------------------------------------------------------------------

The rank-saturation argument above is combinatorial.  We now give a variational reformulation that is arguably deeper.

\begin{theorem}[Free quarks have no non-trivial solution]
\label{thm:confinement_variational}
A configuration containing an isolated A-vertex (not part of any AAA triangle) satisfies $\delta S/\delta\psi = 0$ only trivially, with $S = 0$.
\end{theorem}

\begin{proof}
The Regge action is $S = \sum_h A_h\,\delta_h$ where $A_h = \sqrt{\det(G_h)}$.  A hinge requires 3 vertices.  An isolated A-vertex participates in no complete hinge (all hinges require 3 vertices from the simplex).  Therefore $S = 0$.  With $S = 0$, the dynamical Planck constant $\hbar = A_h/(4\ln 2) = 0$: no spacetime, no time evolution, no observation.

Similarly, a pair of A-vertices (e.g.\ a diquark $ud$) has $\binom{2}{3} = 0$ hinges, so $S = 0$.

The minimum configuration with $S > 0$ is 3 A-vertices forming one AAA hinge: $\det(G_h^{\mathrm{AAA}}) > 0$ (Lemma~\ref{lem:det}).  This is a baryon.
\end{proof}

\begin{remark}[Confinement is not a force --- five proofs of the same fact]
In the standard picture, quarks are ``confined'' by a linearly rising potential.  In DRLT, confinement is not a dynamical effect but a \emph{structural} one, admitting at least five independent formulations:
\begin{enumerate}
\item \textbf{Combinatorial:} $\binom{3}{3} = 1$ --- only one AAA configuration exists per patch, exhausted in a single hop.
\item \textbf{Rank:} The $\CC^3$ spatial Gram matrix has rank $n_S = 3$, saturated by one AAA hinge.  No second independent configuration exists.
\item \textbf{Variational:} A free quark ($< 3$ A-vertices) forms no hinge, so $S = 0$ and $\hbar = 0$.  No non-trivial solution of $\delta S/\delta\psi = 0$ exists.
\item \textbf{Underdetermination:} At $N_\text{eff} = 1$, any AAA configuration is a solution.  There is no ``unconfined'' solution to compare against---the free state is \emph{undefined}, not forbidden.
\item \textbf{Topological:} $\text{BBB} = \binom{n_T}{3} = \binom{2}{3} = 0$.  A purely temporal closure is impossible; spatial vertices are always required.
\end{enumerate}
That five independent arguments yield the same conclusion---from combinatorics, linear algebra, variational calculus, solution theory, and topology---is the strongest evidence that confinement is not an assumption but a theorem.

Confinement need not be \emph{proved}.  It need not even be \emph{stated}.  A free quark is not a solution of $\delta S/\delta\psi = 0$.  It was never a candidate.
\end{remark}

%---------------------------------------------------------------------
\subsubsection{The Navier--Stokes parallel: structural comparison}
\label{sec:navier_stokes_parallel}
%---------------------------------------------------------------------

The Yang--Mills mass gap and the Navier--Stokes blow-up problem share
a common structure, which we first summarise before giving the
rigorous lattice formulation:

\begin{center}
\begin{tabular}{lcc}
\hline
& Yang--Mills & Navier--Stokes \\
\hline
Question & $\Delta > 0$ on $\RR^4$? & No blow-up on $\RR^3$? \\
On the lattice & $\det > 0 \Rightarrow \Delta > 0$ &
  $\det \in (0,1],\;\tr = N$ $\Rightarrow$ bounded \\
Continuum limit & $\det \to 0$, gap vanishes &
  bounds lost, blow-up ``possible'' \\
\hline
\end{tabular}
\end{center}

The rigorous formulation of Navier--Stokes regularity on the lattice
is deferred to Section~\ref{sec:navier_stokes}, after the Yang--Mills
mass gap proof.

%=====================================================================
\subsection{The Mass Gap}
\label{sec:massgap}
%=====================================================================

\begin{lemma}[Non-degeneracy of the AAA hinge]
\label{lem:det}
For any three vertices $\{a, b, c\}$ spanning $\CC^3$, the
restricted Gram matrix $G|_{\mathrm{AAA}}$ is positive definite:
\begin{equation}
\det(G|_{\mathrm{AAA}}) > 0.
\end{equation}
For the orthonormal basis, $\det(G|_{\mathrm{AAA}}) = 1$ exactly.
\end{lemma}

\begin{proof}
The three spatial state vectors $\psi_a, \psi_b, \psi_c \in \CC^3$
span a 3-dimensional space and are therefore linearly independent.
The Gram matrix of linearly independent vectors is positive definite
(since $\mathbf{c}^\dagger G|_h\, \mathbf{c} =
\|\sum_i c_i \psi_i\|^2 > 0$ for all nonzero $\mathbf{c}$).
Positive definiteness implies $\det > 0$.

For the orthonormal case $\langle \psi_i | \psi_j \rangle = \delta_{ij}$,
we have $G|_{\mathrm{AAA}} = I_3$ and $\det = 1$.
\end{proof}

\begin{lemma}[Deficit angle of the AAA hinge]
\label{lem:delta}
Let $K$ be a simplicial complex on $\CCP^4$ with 6 vertices
$\{S_1, S_2, S_3, T_1, T_2, T_3\}$ forming $\partial\Delta^5$, where
$S_i \in \CC^3$ (spatial) and $T_\alpha \in \CC^2$ (temporal).
Since $\dim_{\CC}(\CC^2) = 2$, the three temporal vertices
satisfy $T_3 = \alpha\, T_1 + \beta\, T_2$ with
$|\alpha|^2 + |\beta|^2 = 1$ for some $\alpha, \beta \in \CC$.
The deficit angle at the AAA hinge is:
\begin{equation}
\boxed{\delta_{\mathrm{AAA}} = \pi.}
\end{equation}
This value is independent of $\alpha, \beta$ --- it depends only on $n_T = 2$.
\end{lemma}

\begin{proof}
\textbf{Step 1: Identify the 4-simplices sharing the AAA hinge.}
The AAA hinge $h = \{S_1, S_2, S_3\}$ is contained in those 4-simplices
of $\partial\Delta^5$ that include all three $S$-vertices plus two of
the three $T$-vertices:
\begin{equation}
\sigma_1 = S_1 S_2 S_3 T_1 T_2, \qquad
\sigma_2 = S_1 S_2 S_3 T_1 T_3, \qquad
\sigma_3 = S_1 S_2 S_3 T_2 T_3.
\end{equation}
There are $\binom{3}{2} = 3$ such 4-simplices, each with the correct
$(3S, 2T)$ split per simplex.

\medskip\noindent
\textbf{Step 2: Reduce dihedral angles to temporal angles.}
Within each $\sigma_k$, the hinge $h$ is shared by exactly two
tetrahedra $\tau_a = h \cup \{T_a\}$ and $\tau_b = h \cup \{T_b\}$.
Since $h$ spans $\CC^3$ and $\CC^3 \perp \CC^2$,
the normal to $h$ within $\tau_a$ lies entirely in $\CC^2$
and points in the direction of $T_a$.  The dihedral angle at $h$
equals the Fubini--Study angle between the temporal vertices:
\begin{equation}
\theta_{\sigma_k}(h) = \arccos\bigl|\langle T_a | T_b \rangle\bigr|.
\end{equation}

\medskip\noindent
\textbf{Step 3: Compute using the $n_T = 2$ constraint.}
Choose $T_1, T_2$ as an orthonormal basis of $\CC^2$.
Then $T_3 = \alpha\, T_1 + \beta\, T_2$ with
$|\alpha|^2 + |\beta|^2 = 1$.
Set $|\alpha| = \cos\gamma$ for some $\gamma \in [0, \pi/2]$,
so $|\beta| = \sin\gamma$. The three dihedral angles are:
\begin{align}
\theta_1 &= \arccos|\langle T_1 | T_2 \rangle|
         = \arccos(0) = \frac{\pi}{2}, \\
\theta_2 &= \arccos|\langle T_1 | T_3 \rangle|
         = \arccos(\cos\gamma) = \gamma, \\
\theta_3 &= \arccos|\langle T_2 | T_3 \rangle|
         = \arccos(\sin\gamma) = \frac{\pi}{2} - \gamma.
\end{align}
The last identity uses $\arccos(\sin\gamma) = \pi/2 - \gamma$
for $\gamma \in [0, \pi/2]$.

\medskip\noindent
\textbf{Step 4: Sum and cancel.}
\begin{equation}
\delta_{\mathrm{AAA}} = 2\pi - (\theta_1 + \theta_2 + \theta_3)
= 2\pi - \frac{\pi}{2} - \gamma - \Bigl(\frac{\pi}{2} - \gamma\Bigr)
= 2\pi - \pi = \pi. \qedhere
\end{equation}
\end{proof}

\begin{remark}[Why $\delta = \pi$ is exact]
The cancellation of $\gamma$ in Step~4 is the key: $\theta_2 + \theta_3
= \gamma + (\pi/2 - \gamma) = \pi/2$ for \emph{any} decomposition of
$T_3$ in the temporal basis.  This is a consequence of the normalisation
$|\alpha|^2 + |\beta|^2 = 1$ in a 2-dimensional space.

More precisely, $\delta_{\mathrm{AAA}} = \pi$ follows from a single
algebraic identity:
\begin{equation}
\arccos(\cos\gamma) + \arccos(\sin\gamma) = \frac{\pi}{2},
\qquad \gamma \in [0, \tfrac{\pi}{2}],
\end{equation}
which is the statement that $\sin = \cos(\pi/2 - \cdot)$.  Thus
$\delta = \pi$ is not a numerical coincidence but an identity forced by
$n_T = 2$.
\end{remark}

\begin{remark}[Dimensional dependence]
For general temporal dimension $n_T$, the deficit angle would depend on
the geometry of the temporal sector. The universality
$\delta_{\mathrm{AAA}} = \pi$ (independent of $\alpha, \beta$) is
specific to $n_T = 2$ and would fail for $n_T \geq 3$.  This
singles out the $(3,2)$ split as the unique substrate supporting a
\emph{universal} (geometry-independent) mass gap.
\end{remark}

\begin{theorem}[Yang--Mills Mass Gap]
\label{thm:massgap}
The spectrum of the lattice Hamiltonian $H$
(Definition~\ref{def:massgap}) on the simplicial complex $K$
over $\CCP^4$ has a strictly positive mass gap $\Delta > 0$
in the gauge-singlet sector of $\mathrm{SU}(3)$.
\end{theorem}

\begin{proof}
We establish $\Delta > 0$ in three steps.

\medskip\noindent
\textbf{Step 1: Correlation length $\xi = 1$.}
By Proposition~\ref{prop:confinement}, any gauge-invariant excitation
in the $\mathrm{SU}(3)$ sector has propagation range exactly~1
(Definition~\ref{def:propagation}).  The connected two-point function
satisfies:
\begin{equation}
\langle \mathcal{O}(h_1)\,\mathcal{O}^\dagger(h_2)
\rangle_{\mathrm{conn}} = 0 \qquad
\text{for all AAA hinges } h_1, h_2 \text{ with } |h_1 - h_2| \geq 2.
\end{equation}
The correlator does not merely decay exponentially --- it vanishes
\emph{identically} beyond one lattice hop.  By
Definition~\ref{def:correlation}, the correlation length is
$\xi = 1$ lattice unit.

\medskip\noindent
\textbf{Step 2: Mass gap from correlation length.}
By Definition~\ref{def:massgap} and the spectral representation of
the transfer matrix (Definition~\ref{def:transfer}):
\begin{equation}
\langle \mathcal{O}(0)\,\mathcal{O}^\dagger(r)
\rangle_{\mathrm{conn}} = \sum_{n \geq 1} |c_n|^2\, e^{-(E_n - E_0)\,r},
\end{equation}
where $c_n = \langle n | \mathcal{O}^\dagger | 0 \rangle$.  Since
this vanishes for $r \geq 2$, the dominant exponential must satisfy
$e^{-\Delta \cdot 2} \leq 0$, which is impossible unless $\Delta > 0$
(as the exponential is strictly positive for any finite $\Delta$).

More precisely, the correlator at $r = 1$ is non-zero (the AAA hinge
carries a non-trivial excitation), while at $r = 2$ it vanishes.
This requires $\Delta \geq (\ln C)/1$ for some $C > 1$.

\medskip\noindent
\textbf{Step 3: Explicit lower bound.}
The energy cost of creating a single AAA-hinge excitation above
the vacuum is the Regge action of one AAA hinge:
\begin{equation}
E_1 - E_0 \;\geq\; \inf_{h \in \mathrm{AAA}}
  \frac{S_h}{\hbar_h}
= \inf_{h} \frac{A_h\,\delta_h}{A_h/(4\ln 2)}
= 4\ln 2\, \cdot\, \delta_{\mathrm{AAA}}
= 4\pi\ln 2,
\end{equation}
where we used the area cancellation $S_h/\hbar_h = 4\ln 2\,\delta_h$
(cf.\ Chapter~\ref{app:path_integral}) and
Lemma~\ref{lem:delta} ($\delta_{\mathrm{AAA}} = \pi$).

In Planck units (where we identify one lattice unit with $\ell_P$), the
Regge action of the minimal excitation gives:
\begin{equation}
\boxed{\Delta \;=\; A_h \cdot \delta_{\mathrm{AAA}}
\;=\; \sqrt{\det(G|_{\mathrm{AAA}})} \times \pi \;>\; 0.}
\end{equation}
For the orthonormal (ideal simplex) configuration:
$\det(G|_{\mathrm{AAA}}) = 1$, so $\Delta = \pi$ in Planck units.
\end{proof}

\begin{corollary}
The lightest glueball mass is $m_G = \Delta = \pi\, m_P$ for the
ideal simplex, where $m_P$ is the Planck mass.  In dimensionless
action units, the gap is $4\pi\ln 2 \approx 8.71$.
\end{corollary}

\begin{remark}[Summary of the proof]
The mass gap rests on exactly three algebraic facts, each an identity
rather than a dynamical result:
\begin{enumerate}
\item[(i)] \textbf{Positivity:} $\det(G_{\mathrm{AAA}}) > 0$ --- three
  spanning vectors in $\CC^3$ are linearly independent
  (Lemma~\ref{lem:det}).
\item[(ii)] \textbf{Curvature:} $\delta_{\mathrm{AAA}} = \pi > 0$ ---
  the $n_T = 2$ constraint forces positive deficit angle, independent
  of temporal geometry (Lemma~\ref{lem:delta}).
\item[(iii)] \textbf{Confinement:} $N_{\mathrm{eff}} = \binom{3}{3} = 1$
  --- rank saturation forbids propagation beyond one hop
  (Proposition~\ref{prop:confinement}).
\end{enumerate}
No dynamics, no renormalization, no path integral.  The mass gap is
geometry.
\end{remark}

%=====================================================================
\subsection{No-Go Theorem for Continuous Spacetime}
\label{sec:nogo}
%=====================================================================

\begin{definition}[Continuum limit]
\label{def:continuum}
The continuous spacetime $\mathcal{M}_{\mathrm{cont}} = \RR^4$ is
obtained from the simplicial complex $K_N$ via:
\begin{equation}
\mathcal{M}_{\mathrm{cont}} = \lim_{N \to \infty} K_N,
\end{equation}
subject to $\mathrm{Tr}(G) = N$ and fixed total volume $V$.
In this limit, the lattice spacing $a \to 0$.
\end{definition}

\begin{lemma}[Sub-determinant inequality]
\label{lem:subdet}
Let $G$ be the $n \times n$ Gram matrix of $n$ vectors in $\CC^d$,
and let $G|_S$ be the principal sub-matrix corresponding to a subset
$S \subseteq \{1, \ldots, n\}$ with $|S| = k$.  Then:
\begin{equation}
\det(G|_S) \leq \prod_{i \in S} G_{ii}.
\end{equation}
In particular, for a hinge $h \subset \sigma$:
$\det(G|_h) \leq \det(G|_\sigma)$ when $G$ is positive semidefinite
and all diagonal entries satisfy $G_{ii} \leq 1$.
\end{lemma}

\begin{proof}
The first inequality is Hadamard's inequality applied to the positive
semidefinite sub-matrix $G|_S$: $\det(G|_S) \leq \prod_{i \in S} G_{ii}$.
For the second: since $G \succeq 0$, all principal minors are
non-negative, and Fischer's inequality gives
$\det(G|_\sigma) \geq \det(G|_h) \cdot \det(G|_{\sigma \setminus h})$
for complementary subsets.  Since
$\det(G|_{\sigma \setminus h}) \leq \prod_{i \in \sigma \setminus h}
G_{ii} \leq 1$ (by normalisation $G_{ii} = 1$), we obtain
$\det(G|_h) \leq \det(G|_\sigma) / \det(G|_{\sigma \setminus h})$.
When $G|_{\sigma \setminus h}$ is non-degenerate, this gives
$\det(G|_h) \leq \det(G|_\sigma)$.
\end{proof}

\begin{lemma}[Degeneration of the hinge area]
\label{lem:degen}
In the continuum limit (Definition~\ref{def:continuum}),
$\det(G|_{\mathrm{AAA}}) \to 0$ for every local hinge.
\end{lemma}

\begin{proof}
The total 4-volume of the observable universe is $V < \infty$.  Each
4-simplex $\sigma$ contributes volume
$\mathrm{Vol}(\sigma) = C_d \cdot \det(G|_\sigma)^{1/2}$, where $C_d$
is a dimensional constant.  With $O(N)$ simplices tiling a fixed volume:
\begin{equation}
V = \sum_\sigma \mathrm{Vol}(\sigma)
  = C_d \sum_\sigma \det(G|_\sigma)^{1/2}
  = \mathrm{const}.
\end{equation}
By Jensen's inequality applied to the concave function $x \mapsto x^{1/2}$:
\begin{equation}
\frac{1}{|\{\sigma\}|} \sum_\sigma \det(G|_\sigma)^{1/2}
\leq \Bigl(\frac{1}{|\{\sigma\}|} \sum_\sigma \det(G|_\sigma)\Bigr)^{1/2}.
\end{equation}
Since the left-hand side equals $V / (C_d \cdot |\{\sigma\}|)$
and $|\{\sigma\}| = O(N)$:
\begin{equation}
\langle \det(G|_\sigma) \rangle \geq \frac{V^2}{C_d^2 \cdot N^2}
\to 0 \quad \text{as } N \to \infty.
\end{equation}
More directly: $\langle \det^{1/2} \rangle = V / (C_d N) \to 0$.
Since $\det^{1/2} \geq 0$, the average tending to zero forces
$\det(G|_\sigma) \to 0$ for a.e.\ simplex.

By Lemma~\ref{lem:subdet}: $\det(G|_{\mathrm{AAA}}) \leq
\det(G|_\sigma) \to 0$ for each hinge $h \subset \sigma$.
\end{proof}

\begin{lemma}[Collapse of the gauge fiber]
\label{lem:fiber}
If $\det(G|_h) \to 0$ for all hinges $h$, all holonomies become
trivial.
\end{lemma}

\begin{proof}
The holonomy around a triangle $h = \{i, j, k\}$ is:
\begin{equation}
\Phi_h = \arg(G_{ij}\, G_{jk}\, G_{ki}).
\end{equation}
The $3 \times 3$ hinge determinant expands as (Chapter~\ref{ch:simplex_geometry}):
\begin{equation}
\det(G|_h) = 1 - |G_{ij}|^2 - |G_{jk}|^2 - |G_{ki}|^2
  + 2\,|G_{ij}|\,|G_{jk}|\,|G_{ki}|\cos\Phi_h,
\end{equation}
where we used $G_{ii} = 1$.  Rearranging:
\begin{equation}
\cos\Phi_h = \frac{\det(G|_h) - 1 + |G_{ij}|^2 + |G_{jk}|^2 + |G_{ki}|^2}
{2\,|G_{ij}|\,|G_{jk}|\,|G_{ki}|}.
\label{eq:cosPhidet}
\end{equation}

As $\det(G|_h) \to 0$ with fixed volume, neighbouring vertices approach
each other: $|G_{ij}| \to 1$.  Substituting
$|G_{ij}| = |G_{jk}| = |G_{ki}| = 1 - \epsilon$ with
$\epsilon \to 0$:
\begin{equation}
\cos\Phi_h \to \frac{0 - 1 + 3(1 - 2\epsilon)}{2(1 - 3\epsilon)}
= \frac{2 - 6\epsilon}{2 - 6\epsilon} \to 1.
\end{equation}
Therefore $\Phi_h \to 0$ for every hinge.  All holonomies become
trivial, and by the Ambrose--Singer theorem \cite{Ambrose-Singer},
a connection with trivial holonomy is flat: there is no Yang--Mills
dynamics.
\end{proof}

\begin{remark}[Ambrose--Singer and the discrete setting]
The Ambrose--Singer theorem states that the Lie algebra of the holonomy
group is generated by the curvature 2-forms.  In the discrete setting,
the ``curvature'' at hinge $h$ is $\Phi_h$.  When all $\Phi_h \to 0$,
the holonomy group collapses to the identity.  This is not a technical
failure of a particular regularisation scheme but the geometric statement
that gauge flux requires a finite area to enclose: shrinking all areas
to zero destroys all non-trivial gauge structure.
\end{remark}

\begin{theorem}[No-Go]
\label{thm:nogo}
A strictly positive Yang--Mills mass gap ($\Delta > 0$) and a continuous
spacetime ($\RR^4$) satisfying the Wightman axioms are mutually
exclusive.
\end{theorem}

\begin{proof}
By Theorem~\ref{thm:massgap}, the mass gap on the lattice is:
\begin{equation}
\Delta = \sqrt{\det(G|_{\mathrm{AAA}})} \times \delta_{\mathrm{AAA}}
= \sqrt{\det(G|_{\mathrm{AAA}})} \times \pi.
\end{equation}

The Wightman axioms require the quantum field theory to be formulated on
$\RR^4$, which implies the continuum limit
(Definition~\ref{def:continuum}).  We show the limit destroys the gap
by two independent mechanisms.

\medskip\noindent
\textbf{(a) The mass gap vanishes.}
By Lemma~\ref{lem:degen}, $\det(G|_{\mathrm{AAA}}) \to 0$ in the
continuum limit.  Since $\delta_{\mathrm{AAA}} = \pi$ is a fixed
constant (Lemma~\ref{lem:delta}):
\begin{equation}
\lim_{N \to \infty} \Delta
= \pi \cdot \lim_{N \to \infty} \sqrt{\det(G|_{\mathrm{AAA}})}
= \pi \cdot 0 = 0.
\end{equation}

\medskip\noindent
\textbf{(b) The gauge structure collapses.}
By Lemma~\ref{lem:fiber}, all $\mathrm{SU}(3)$ holonomies become
trivial ($\Phi_h \to 0$).  The confinement mechanism
(Proposition~\ref{prop:confinement}) becomes vacuous: there is no
non-trivial gauge field left to confine.

\medskip\noindent
The mass gap vanishes for two independent and logically sufficient
reasons: (a)~its numerical value goes to zero, and (b)~the gauge
symmetry that defines it ceases to exist.
\end{proof}

\begin{corollary}
The Clay Millennium Problem, as formulated on $\RR^4$ with the
Wightman axioms, has no solution with $\Delta > 0$. The mass gap exists
if and only if spacetime is discrete.
\end{corollary}

\begin{remark}[Why renormalization cannot help]
In standard lattice gauge theory, the bare coupling $g_0(a)$ is tuned as
$a \to 0$ so that physical quantities $m_{\mathrm{phys}} = m(a)/a$
remain finite. This procedure \emph{assumes} that the lattice is a
computational approximation to a continuum theory.

In DRLT, there is no bare coupling to tune. The coupling
$\agut = 6/(25\pi^2)$ is a mathematical constant
derived from the geometry of $\CC^5$ . The hinge
area $A_h = \sqrt{\det(G|_h)}$ is a physical observable (the local
Planck constant $\hbar_h = A_h/(4\ln 2)$), not a regularization
artifact. ``Renormalizing'' $A_h$ would mean changing the local laws of
physics. There is no freedom to absorb $\det \to 0$ into a redefinition
of parameters: the lattice is not an approximation to be refined; it
\emph{is} the physics.
\end{remark}

%=====================================================================
\subsection{Navier--Stokes Regularity on the Lattice}
\label{sec:navier_stokes}
%=====================================================================

We now formalise the structural parallel between the Yang--Mills mass
gap and the Navier--Stokes regularity problem.  The Clay formulation
of the Navier--Stokes problem \cite{Clay} asks:

\begin{quote}
\emph{For smooth, divergence-free initial data on $\RR^3$ (or
$\mathbb{T}^3$) with finite energy, do smooth solutions of the
3D incompressible Navier--Stokes equations exist for all time?}
\end{quote}

\noindent
We prove that on the DRLT lattice, the answer is unconditionally yes:
all discrete analogues of the Sobolev norms remain bounded for all time.
We then show that the continuum limit destroys these bounds by the same
mechanism that destroys the mass gap.

%---------------------------------------------------------------------
\subsubsection{The lattice fluid model}
%---------------------------------------------------------------------

\begin{definition}[Lattice velocity field]
\label{def:lattice_velocity}
On the simplicial complex $K$ with $N$ vertices, the
\emph{velocity field} is the discrete gradient of the geometric weight:
\begin{equation}
v_{ij} = W_j - W_i, \qquad
W_i = \frac{1}{d} \sum_{j \neq i} |G_{ij}|^2,
\end{equation}
where the sum runs over vertices $j$ adjacent to $i$ in $K$.  The
velocity $v_{ij}$ is defined on each edge $(i,j)$ of $K$.
\end{definition}

\begin{definition}[Lattice vorticity and higher derivatives]
\label{def:lattice_vorticity}
The \emph{vorticity} at hinge $h = \{i, j, k\}$ is the discrete curl:
\begin{equation}
\omega_h = v_{ij} + v_{jk} + v_{ki} = 0
\end{equation}
(identically zero by telescoping, since $v_{ij} = W_j - W_i$).
The $m$-th order \emph{discrete derivative} on $K$ is:
\begin{equation}
(\nabla^m v)_{i} = \sum_{\substack{j_1, \ldots, j_m \\ \text{path from } i}}
\prod_{l=1}^{m} v_{j_{l-1}\, j_l},
\end{equation}
summing over paths of length $m$ starting at $i$.
\end{definition}

\begin{definition}[Discrete Sobolev norm]
\label{def:sobolev}
The discrete $H^s$ norm of the velocity field $v$ on $K$ is:
\begin{equation}
\|v\|_{H^s(K)}^2 = \sum_{m=0}^{s}
  \sum_{\text{edges } (i,j)} |(\nabla^m v)_{ij}|^2.
\end{equation}
\end{definition}

%---------------------------------------------------------------------
\subsubsection{Regularity on the finite lattice}
%---------------------------------------------------------------------

\begin{lemma}[Pointwise velocity bound]
\label{lem:velocity_bound}
For all edges $(i,j)$ in $K$:
\begin{equation}
|v_{ij}| \leq c = 2.
\end{equation}
\end{lemma}

\begin{proof}
By Cauchy--Schwarz, $|G_{ij}|^2 = |\langle\psi_i|\psi_j\rangle|^2
\leq \|\psi_i\|^2 \|\psi_j\|^2 = 1$ (since $G_{ii} = 1$).
Therefore $W_i = \frac{1}{d}\sum_j |G_{ij}|^2 \in [0, 1]$ for each
vertex.  Hence $|v_{ij}| = |W_j - W_i| \leq 1$.

More precisely, using $W_i \in [0, (N-1)/d]$ and the triangle
inequality on the lattice graph where each vertex has at most
$O(d)$ neighbours in a single simplex:
$|v_{ij}| \leq 2$ (the lattice speed of light $c = 2$ in the
DRLT framework).
\end{proof}

\begin{lemma}[Determinant and trace bounds]
\label{lem:det_trace_bounds}
For all hinges $h$ in $K$ and at all times:
\begin{enumerate}
\item[(i)] $\det(G|_h) \in [0, 1]$ \quad (Hadamard inequality with
  $G_{ii} = 1$),
\item[(ii)] $\tr(G) = N$ \quad (normalisation conservation),
\item[(iii)] all eigenvalues of $G$ lie in $[0, N]$ \quad
  ($G \succeq 0$ and $\tr(G) = N$).
\end{enumerate}
\end{lemma}

\begin{proof}
(i) By Hadamard's inequality, $\det(G|_h) \leq \prod_{i \in h} G_{ii}
= 1$ for any principal minor.  Non-negativity follows from
$G \succeq 0$.

(ii) By construction, $\tr(G) = \sum_i \|\psi_i\|^2 = N$.

(iii) The eigenvalues $\lambda_k$ satisfy $\lambda_k \geq 0$
(positive semidefiniteness) and $\sum_k \lambda_k = N$ (trace).
Hence $\lambda_k \leq N$ for all $k$.
\end{proof}

\begin{theorem}[Lattice regularity --- no blow-up]
\label{thm:lattice_regularity}
Let $K$ be a finite simplicial complex with $N$ vertices on $\CCP^4$.
Let $v(t)$ be the velocity field (Definition~\ref{def:lattice_velocity})
evolving under the Regge action $S = \sum_h A_h\,\delta_h$.
Then for all $s \geq 0$ and all $t > 0$:
\begin{equation}
\boxed{\|v(t)\|_{H^s(K)} \;\leq\; C(N, s) \;<\; \infty,}
\end{equation}
where $C(N,s) = 2^s \cdot N^{(s+1)/2} \cdot |\{E\}|^{1/2}$ and
$|\{E\}|$ is the number of edges in $K$.
\end{theorem}

\begin{proof}
\textbf{Step 1: Pointwise bound on derivatives.}
By Lemma~\ref{lem:velocity_bound}, $|v_{ij}| \leq 2$ for each edge.
Since $\nabla^m v$ is a sum of products of $m$ velocity values along
paths of length $m$, and each vertex in $K$ has at most $D$ neighbours
(where $D$ is the maximum vertex degree):
\begin{equation}
|(\nabla^m v)_{ij}| \leq D^{m-1} \cdot 2^m.
\end{equation}

\textbf{Step 2: Sum over edges.}
The complex $K$ has at most $\binom{N}{2}$ edges (finite).  Therefore:
\begin{equation}
\|v\|_{H^s(K)}^2 = \sum_{m=0}^{s} \sum_{(i,j)} |(\nabla^m v)_{ij}|^2
\leq (s+1) \cdot \binom{N}{2} \cdot D^{2(s-1)} \cdot 4^s < \infty.
\end{equation}

\textbf{Step 3: Time independence of the bound.}
The bounds in Lemma~\ref{lem:det_trace_bounds} are
\emph{algebraic constraints} ($G \succeq 0$, $\tr = N$, $G_{ii} = 1$),
not dynamical ones.  They hold at $t = 0$ by construction and are
preserved by any unitary or gradient-flow evolution that maintains the
Gram matrix structure.  Therefore $C(N,s)$ does not depend on $t$.
\end{proof}

\begin{remark}[Why the theorem is trivial --- and why that is the point]
The proof is elementary: finite sums of bounded quantities are finite.
This is precisely the point.  On a \emph{finite} simplicial complex,
blow-up is algebraically impossible.  The Navier--Stokes blow-up
problem is hard because it asks about $\RR^3$, where sums become
integrals and pointwise bounds are lost.
\end{remark}

%---------------------------------------------------------------------
\subsubsection{The continuum limit destroys regularity}
%---------------------------------------------------------------------

\begin{theorem}[No-go for Navier--Stokes regularity via continuum limit]
\label{thm:ns_nogo}
In the continuum limit $N \to \infty$ (Definition~\ref{def:continuum}),
the uniform Sobolev bound of Theorem~\ref{thm:lattice_regularity}
diverges: $C(N, s) \to \infty$ for all $s \geq 1$.
\end{theorem}

\begin{proof}
The bound $C(N, s) = 2^s \cdot N^{(s+1)/2} \cdot |\{E\}|^{1/2}$
grows as $N^{(s+1)/2} \cdot N \to \infty$ (since
$|\{E\}| = O(N)$ for a simplicial complex of bounded degree).
For $s \geq 1$, $C(N,s) \to \infty$ as $N \to \infty$.

The blow-up protection is lost because:
\begin{enumerate}
\item \textbf{Finite sums $\to$ infinite series}: with $N \to \infty$
  edges, the discrete Sobolev norm becomes an infinite sum that
  can diverge.
\item \textbf{Lattice spacing $\to$ 0}: as $a \to 0$, the velocity
  gradient $v_{ij}/a$ can grow without bound even if $v_{ij}$ is
  bounded, because the denominator vanishes.
\item \textbf{Pointwise bounds dilute}: $W_i = \frac{1}{d}\sum_j
  |G_{ij}|^2$ sums over $O(N)$ terms in the continuum limit.
  While each $|G_{ij}|^2 \leq 1$, the sum can reach $O(N/d)$,
  making $v_{ij} = O(N/d)$ unbounded.
\end{enumerate}
\end{proof}

\begin{theorem}[Structural equivalence of the two millennium problems]
\label{thm:structural_equivalence}
The Yang--Mills mass gap and the Navier--Stokes regularity problems
have the same logical structure on the DRLT lattice:
\begin{enumerate}
\item Both properties hold unconditionally on any finite simplicial
  complex $K$ with $N < \infty$
  (Theorems~\ref{thm:massgap} and~\ref{thm:lattice_regularity}).
\item Both properties are destroyed by the continuum limit
  $N \to \infty$
  (Theorems~\ref{thm:nogo} and~\ref{thm:ns_nogo}).
\item In both cases, the mechanism is the same:
  \begin{itemize}
  \item Mass gap: $\det(G|_{\mathrm{AAA}}) \to 0 \Rightarrow
    \Delta \to 0$.
  \item Regularity: $C(N,s) \to \infty \Rightarrow$ blow-up possible.
  \end{itemize}
\end{enumerate}
\end{theorem}

\begin{proof}
This is a direct consequence of Theorems~\ref{thm:massgap},
\ref{thm:nogo}, \ref{thm:lattice_regularity},
and~\ref{thm:ns_nogo}.
\end{proof}

\begin{remark}[Two millennium problems, one answer]
Both the Yang--Mills mass gap and the Navier--Stokes regularity
problem ask: ``Prove property $X$ on $\RR^n$.''  Both are open
because $\RR^n$ lacks the algebraic bounds ($\det \in (0,1]$,
$\tr = N$, finite sums) that the discrete structure provides for free.
The resolution is not a cleverer proof on $\RR^n$ but the recognition
that $\RR^n$ is an approximation that does not preserve the relevant
bounds.

The minimum viscosity $\eta/s = 1/(4\pi)$
(Appendix~\ref{app:qcd}) is the lattice's intrinsic flow resistance
--- the geometric lower bound on dissipation that makes blow-up
impossible.
\end{remark}

%=====================================================================
\subsection{Discussion}
\label{sec:discussion}
%=====================================================================

\subsubsection{Lattice QCD versus DRLT}

The distinction is ontological:

\begin{center}
\begin{tabular}{lcc}
\hline
& Lattice QCD & DRLT \\
\hline
Lattice is & computational tool & physical reality \\
Continuum limit & required & forbidden \\
Mass gap & observed numerically & proved algebraically \\
Coupling & bare $g_0(a)$, tuned & $\alpha = 6/(25\pi^2)$, derived \\
Confinement & from flux tubes & from $\dim(\CC^3) = 3$ \\
\hline
\end{tabular}
\end{center}

\noindent
Lattice QCD and DRLT agree on all observable predictions (confinement,
mass gap, asymptotic freedom), but disagree on what the lattice \emph{is}.
DRLT claims the lattice is fundamental; lattice QCD treats it as an
approximation. Our results show that the mass gap \emph{requires}
the lattice to be fundamental.

\subsubsection{The Millennium Problem reconsidered}

The Clay problem asks: ``Prove $\Delta > 0$ on $\RR^4$.''
Theorem~\ref{thm:nogo} shows that $\Delta > 0$ and $\RR^4$ are
mutually exclusive. This does not mean the mass gap is false --- it
means the premise ($\RR^4$) is wrong. The correct formulation:

\begin{quote}
\emph{The Yang--Mills mass gap exists on the natural simplicial complex
$K$ over $\CCP^4$, with $\Delta = \sqrt{\det(G_{\mathrm{AAA}})}
\times \pi$, and it cannot exist on $\RR^4$.}
\end{quote}

\subsubsection{Numerical confirmation}

The algebraic proof is self-contained, but the key quantities have been
confirmed numerically in DRLT experiments:

\begin{itemize}
\item $\det(G|_{\mathrm{AAA}}) = 1.000000 \pm 0$ for orthogonal
  $\CC^3$ basis (EXP-043b).
\item Deficit angle $\delta_{\mathrm{AAA}} = 180.00^\circ$ (variational, EXP-047b)
  (EXP-042a).
\item Strong-sector correlation: $W_S$ local/global ratio $> 2$,
  confirming confinement (EXP-029).
\item $\det(G_h) > 0$ for all physical configurations tested across
  $N = 10$--$10{,}000$ vertices (EXP-032, EXP-043).
\end{itemize}

\subsubsection{Why the problem was hard}

The Yang--Mills mass gap problem resisted proof for decades because it
was formulated in the wrong language. Continuous spacetime
($\RR^4$) destroys the very geometric structure that creates the
mass gap. The problem is analogous to proving the discreteness of atomic
spectra while insisting on classical mechanics: the framework excludes
the answer.

The resolution is not a more clever analytic technique on $\RR^4$.
It is the recognition that spacetime is a simplicial complex on
$\CCP^4$, that the strong force lives on AAA hinges, and that
$\binom{3}{3} = 1$ is the entire content of confinement.
